%%%PAGE 191 rot=0 legibility=good summary=卫星面积速率；中子；缆绳卫星离轨能量；氢原子巴尔末系；天体质量变化(月球土运回地球)；射线来源；天问一号；月球着陆反冲力；两椭圆/圆轨道能量比较(含图)
%%%BLOCK id=191-01 ch=05 sec=卫星运行规律·面积速率 kind=problem alt=none cont=none y=0.04-0.19
% 页面右上角叠放着另一张纸(化学笔记)的边缘，其上文字(“n CO2 … 燃烧得CO和CO2混合气体 … CO为0.4·3/4×22.4=6.72L”等)属相邻页，未转录
1.
\begin{problem}
卫星绕地球匀圆，地球半径$R$，表面重力加速度$g$，轨道半径$r$，则卫星与地心连线单位时间扫过的面积？ % 原稿“表面”二字为小字，插写在“重力”上方
\begin{solution}
\[
\frac{GMm}{r^2}=\frac{mv^2}{r}\qquad \frac{GMm'}{R^2}=m'g\qquad S=\frac{vr}{2}=\frac{R}{2}\sqrt{gr}
\]
\end{solution}
\end{problem}
%%%END
%%%BLOCK id=191-02 ch=17 sec=原子核·中子 kind=knowledge alt=none cont=none y=0.165-0.235
2.\ 用电磁场来加速减速来\unclear{束缚}中子\ $\times$\qquad 中子不带电不受电磁场作用 % 原稿“来”后一字难辨(缚/束缚?)
%%%END
%%%BLOCK id=191-03 ch=05 sec=卫星离轨·能量 kind=knowledge alt=11,12 cont=none y=0.19-0.32
3.
%FIG p191_a 0.07 0.19 0.275 0.31
\notefig[0.25]{figures/p191_a.png}

\underline{安培力处处不等（地磁不均匀）} % 原稿此行下划线，并被一条弧线(自图中“空心导体”下方向右上延伸)斜穿

离轨\quad 安培力做负功，万有引力做正功\quad 故动能不一定减小。
%%%END
%%%BLOCK id=191-04 ch=17 sec=氢原子能级与光谱 kind=problem alt=none cont=none y=0.30-0.46
4.
\begin{problem}
$E_n=\dfrac{E_1}{n^2}$\quad $n=2,3,4$。H从$n=2\to n=1$放出光子频率$\nu$，巴尔末系指H从高能级向$n=2$能级跃迁时释放的光子，使

使处于基态的氢原子释放出巴尔末系的光子的照射光光子最小频率？ % CHECK: 上一行末尾的“使”与本行开头的“使”重复，照原样转录
\begin{solution}
\[
\frac{E_1}{4}-E_1=h\nu
\]
\[
\frac{E_1}{9}-E_1=h\nu'\qquad \nu'=\frac{32}{27}\nu
\]
\end{solution}
\end{problem}
%%%END
%%%BLOCK id=191-05 ch=05 sec=天体质量变化·运行参量 kind=note alt=none cont=none y=0.44-0.51
5.\ 若$M_{\text{地球}}\uparrow$，$M_{\text{月球}}$不变，\ $F_{\text{引}}=\dfrac{GMm}{R^2}\uparrow$\quad $v=\sqrt{\dfrac{GM}{R}}\uparrow$\quad $T=\dfrac{2\pi R}{v}\downarrow$，$R\downarrow$\ （月球做近心运动） % F的下标原稿形似“31”，按“引”转录
%%%END
%%%BLOCK id=191-06 ch=17 sec=核反应·守恒规律 kind=knowledge alt=none cont=none y=0.485-0.52
6.\ 核反应遵循\underline{电荷守恒}\ \underline{质量数守恒}\ 但\underline{质量不守恒}
%%%END
%%%BLOCK id=191-07 ch=05 sec=天体质量密度·环绕天体 kind=note alt=none cont=none y=0.52-0.58
7.\ 无法算环绕天体质量（及力、动量、功能）与$m$有关） % CHECK: 括号不配对(原稿共一个左括号、两个右括号)，照原样转录
%%%END
%%%BLOCK id=191-08 ch=17 sec=原子能级跃迁·射线来源 kind=knowledge alt=none cont=none y=0.55-0.65
8.\ 氢原子从高能级向低能级跃迁不产生$\gamma$射线，（$\gamma$射线是原子核能级跃迁时产生的）

\hspace{1.5em}从$n=4$向低能级跃迁时辐射出6种频率的光，只有2种可见光。

\hspace{1.5em}红外线\ 可见光\ 紫外线是原子外层电子受激发产生的

\hspace{1.5em}$X$射线是原子内层电子受激发产生的。$\gamma$射线是原子核的能级跃迁产生的
%%%END
%%%BLOCK id=191-09 ch=05 sec=变轨·天问一号 kind=note alt=none cont=none y=0.65-0.72
9.\ 天问一号到火星轨道\unclear{时}与地球间距\underline{不是}极值（不同$T$） % CHECK: “轨道”后一字(及其上方的小字)字迹缠绕难辨，按语义猜“时”

匀圆\quad $\vec{a}\ \vec{v}\ \vec{F}\ \vec{p}$都变化，$E_k$不变
%%%END
%%%BLOCK id=191-10 ch=05 sec=月球着陆·反冲力 kind=note alt=02,03 cont=none y=0.72-0.80
10.\ $g_{\text{月}}=1.7\unit{m/s^2}$\quad $v=\sqrt{2g_{\text{月}}h}$\quad 飞船
%FIG p191_b 0.431 0.733 0.515 0.785
\notefig[0.25]{figures/p191_b.png}

$F_{\text{反冲力}}=mg_{\text{月}}$\quad 飞船着陆过程$E_{\text{机}}$不守恒
%%%END
%%%BLOCK id=191-11 ch=05 sec=天体质量变化·运行参量 kind=note alt=none cont=none y=0.80-0.87
11.\ 把月\unclear{地}土运\unclear{回}地球，中心距离不变\quad $F=\dfrac{GMm}{r^2}\downarrow$\quad $M+m=C$\quad $M-m\uparrow$\quad $Mm\downarrow$ % CHECK: “月”后一字、“运”后一字难辨；意思应为把月球上的土运回地球

\hspace{4em}$T=2\pi\sqrt{\dfrac{r^3}{GM}}\downarrow$\quad $M\uparrow$\quad $a_{\text{向}}=\dfrac{GM}{r^2}\uparrow$\quad $g=\dfrac43\pi\rho GR$\quad $R\downarrow\ g\downarrow$
%%%END
%%%BLOCK id=191-12 ch=17 sec=核反应·人工转变 kind=knowledge alt=none cont=none y=0.865-0.89
卢瑟福发现质子\quad 查德威克发现中子\qquad 慢中子与\unclear{${}^{6}$Li}反应\ \underline{是人工转变}\ \underline{不是裂变反应} % 原稿“Li”左上角的数字形似“5”，按6猜
%%%END
%%%BLOCK id=191-13 ch=05 sec=卫星变轨·椭圆轨道能量 kind=derivation alt=none cont=none y=0.885-1.00
12.
%FIG p191_c 0.055 0.886 0.285 0.957
\notefig[0.25]{figures/p191_c.png}

2\unclear{的椭圆}长轴为$a$

两卫星到C时速率相等\quad $-\dfrac{GMm}{2a}=-\dfrac{GMm}{2R}$\quad $\therefore R=a$\quad 故$\unclear{OA}=\unclear{OC}=R=a$

势能相同\qquad $E_{\text{机}}$也相同

由$\dfrac{R^3}{T_1^2}=\dfrac{a^3}{T_2^2}$\quad $\therefore T_1=T_2$\qquad $E_{\text{机}}=-\dfrac{GMm}{2a}$\qquad 近地$v_3>v_2$\qquad $E_2<E_3$
%%%END
%%%PAGEEND 191
